\label{chap:mps}

\emph{Standing assumption.}  Unless stated otherwise, $(X,\Abb,m,T)$ is an invertible measure
preserving system in which $X$ is a compact metric space and $\Abb$ is its Borel
$\sigma$-algebra.  Proposition~\ref{prop:reduction} removes this assumption at the end.

\section{Systems, factors, spectral measures}

\begin{definition}[Measure preserving system]
  \label{def:mps}
  \lean{CantorSzemeredi.MPS}
  \leanok
  A probability space $(X,\Abb,m)$ together with a measurable $T : X \to X$ with
  $m(T^{-1}A) = m(A)$ for all $A \in \Abb$.  \emph{Invertible} means $T$ is a measurable
  bijection with measurable inverse.  \emph{Ergodic} means every $T$-invariant set has measure
  $0$ or $1$.
\end{definition}

\begin{remark}[Lean design note]
  Mathlib has \texttt{MeasureTheory.MeasurePreserving} and
  \texttt{MeasureTheory.Ergodic}; the Koopman operator is available as composition with $T$ on
  $L^p$.  A light structure bundling $(X,\Abb,m,T)$ is worth introducing because the whole
  development quantifies over systems and over powers $T^k$ of a single system.
\end{remark}

\begin{definition}[Spectral measure]
  \label{def:spectral-measure}
  \lean{CantorSzemeredi.HasSpectralMeasure, CantorSzemeredi.Torus,
    CantorSzemeredi.HasSpectralMeasure.norm_sum_sq,
    CantorSzemeredi.HasSpectralMeasure.inner_zpow,
    CantorSzemeredi.HasSpectralMeasure.measure_univ_eq,
    CantorSzemeredi.hasSpectralMeasure_specMeasure, CantorSzemeredi.specMeasure}
  \leanok
  For a unitary $U$ on a Hilbert space and a vector $v$, $\sigma_v$ is the finite positive Borel
  measure on $\TT = \RR/\ZZ$ with $\langle U^n v, v \rangle = \int_{\TT} e(n\alpha) \, d\sigma_v(\alpha)$
  for all $n \in \ZZ$.  When no operator is named, $U$ is the Koopman operator $f \mapsto f \circ T$.
\end{definition}

\begin{remark}[How the formalization is arranged, and the route to existence]
  \leanok
  Herglotz/Bochner for positive definite sequences on $\ZZ$ is not in Mathlib, and neither is
  the weak-$\ast$ compactness of finite measures that the classical Fej\'er kernel proof uses.
  The formalization is therefore arranged in two layers, in
  \texttt{Dynamics/Spectral.lean} and \texttt{Dynamics/SpectralExists.lean}.

  \emph{The interface.}  \texttt{HasSpectralMeasure U v \(\sigma\)} is a predicate recording the
  defining property, and all consequences are proved from it.  Chief among them is
  \texttt{norm\_sum\_sq},
  \[
    \Bigl\lVert \sum_i a_i U ^{k_i} v \Bigr\rVert ^2
      = \int_{\TT} \Bigl\lvert \sum_i a_i e(k_i \alpha) \Bigr\rvert ^2 \;d \sigma ,
  \]
  the form in which the spectral measure is used below: with $a$ the constants of an average
  and $k_i = s k_n$, the left side is the $L^2$ norm of an ergodic average and the right side
  is what Lemma~\ref{lem:irrational} estimates.

  \emph{Existence.}  \texttt{hasSpectralMeasure\_specMeasure} constructs $\sigma_v$ for any
  unitary on a complex Hilbert space, by the Riesz--Markov--Kakutani route: the continuous
  functional calculus makes $f \mapsto \operatorname{Re}\langle v, f(U)v\rangle$ a positive
  linear functional on $C(\TT,\RR)$, and \texttt{RealRMK.rieszMeasure} represents it.  Two
  points of the construction are worth recording.
  \begin{itemize}
    \item Mathlib's \texttt{cfc\_nonneg} is stated for an \emph{ordered} scalar ring, and $\CC$
      carries no \texttt{LE} instance, so it is unavailable for the $\CC$-valued calculus of a
      normal operator.  Positivity is obtained instead by factoring: a non-negative $f$ on the
      spectrum is $\overline g g$ for $g = \sqrt f$, and $x^\ast x \geq 0$ in any
      C$^\ast$-algebra.
    \item \texttt{cfc} consumes functions $\CC \to \CC$, which would force extending a function
      on $\TT$ off the circle.  Using \texttt{cfcHom}, which consumes
      $C(\operatorname{spec} a, \CC)$ directly, avoids this.
  \end{itemize}
  The Fourier coefficients are then identified by splitting $e(n\alpha)$ into real and
  imaginary parts and recombining: the two pullbacks add up to $z \mapsto z^n$ on the spectrum,
  which \texttt{cfcHom} carries to $U^n$ (\texttt{cfcHom\_zpowSpec}).
\end{remark}

\begin{lemma}[Wiener's lemma]
  \label{lem:wiener}
  \lean{CantorSzemeredi.cesaroChar, CantorSzemeredi.tendsto_cesaroChar_of_ne_zero,
    CantorSzemeredi.norm_cesaroChar_le_one, CantorSzemeredi.cesaroChar_zero}
  For a finite positive measure $\sigma$ on $\TT$,
  $\sum_{\theta} \sigma(\{\theta\})^2 = \lim_N \avg{1 \leq n \leq N} \abs{\widehat\sigma(n)}^2$.
  In particular $\sigma$ is atomless iff $\avg{1 \leq n \leq N} \abs{\widehat\sigma(n)}^2 \to 0$.
\end{lemma}

\begin{definition}[Weak mixing]
  \label{def:weak-mixing}
  \lean{CantorSzemeredi.MPS.WeaklyMixing}
  \leanok
  \uses{def:mps}
  $(X,\Abb,m,T)$ is \emph{weakly mixing} if
  $\lim_N \avg{1 \leq n \leq N} \abs{m(A \cap T^{-n}B) - m(A)m(B)} = 0$ for all $A,B \in \Abb$.
  Equivalently, every mean-zero $f \in L^2(X)$ has atomless spectral measure.
\end{definition}

\begin{definition}[Compact system]
  \label{def:compact-system}
  \lean{CantorSzemeredi.MPS.CompactSystem, CantorSzemeredi.MPS.CompactFunction}
  \leanok
  \uses{def:mps}
  $f \in L^2(X)$ is \emph{compact} if $\overline{\{T^n f : n \in \ZZ\}}$ is compact in $L^2(X)$;
  the system is \emph{compact} if every $f \in L^2(X)$ is compact.
\end{definition}

\begin{definition}[Factor, extension]
  \label{def:factor}
  \lean{CantorSzemeredi.MPS.Factor, CantorSzemeredi.MPS.Factor.trivial}
  \leanok
  \uses{def:mps}
  A \emph{factor} of $(X,\Abb,m,T)$ is a $T$-invariant sub-$\sigma$-algebra $\Dbb \subset \Abb$,
  factors agreeing up to null sets being identified.  If $\Dbb \subset \Ebb \subset \Abb$ are
  factors, $\Ebb$ is an \emph{extension} of $\Dbb$.
\end{definition}

\section{Conditional structures}

\begin{definition}[Conditional inner product, Hilbert module]
  \label{def:hilbert-module}
  \lean{CantorSzemeredi.MPS.condInner, CantorSzemeredi.MPS.condNorm}
  \leanok
  \uses{def:factor}
  For $f,g \in L^2(\Ebb)$ put $\cip{f}{g} \coloneqq \Exp_m(f \bar g \mid \Dbb) \in L^2(\Dbb)$,
  with associated conditional norm $\cnorm{f}$.  The \emph{Hilbert module}
  $L^2(\Ebb \mid \Dbb)$ consists of the $f \in L^2(\Ebb)$ with $\cnorm{f} \in L^\infty(\Dbb)$.
\end{definition}

\begin{lemma}[Conditional Cauchy--Schwarz and triangle inequalities]
  \label{lem:cond-cs}
  \lean{CantorSzemeredi.MPS.abs_condInner_le,
    CantorSzemeredi.MPS.condInner_quadratic_nonneg}
  \leanok
  \uses{def:hilbert-module}
  $\abs{\cip{f}{g}} \leq \cnorm{f}\,\cnorm{g}$ and $\cnorm{f+g} \leq \cnorm{f} + \cnorm{g}$,
  almost everywhere.
\end{lemma}

\begin{lemma}[Conditional expectation commutes with $T$]
  \label{lem:condexp-commutes}
  \lean{CantorSzemeredi.MPS.condExp_comp}
  \leanok
  \uses{def:factor}
  If $\Dbb$ is $T$-invariant then $\Exp(a \circ T \mid \Dbb) = \Exp(a \mid \Dbb) \circ T$.
\end{lemma}

\begin{definition}[Conditionally weakly mixing function; weakly mixing extension]
  \label{def:cwm}
  \lean{CantorSzemeredi.MPS.CondWeaklyMixing, CantorSzemeredi.MPS.WeaklyMixingExtension}
  \leanok
  \uses{def:hilbert-module}
  $f \in L^2(\Ebb)$ is \emph{conditionally weakly mixing} if
  $\lim_N \avg{1 \leq n \leq N} \abs{\cip{T^n f}{f}} = 0$ in $L^2(\Dbb)$.
  The extension $\Ebb \mid \Dbb$ is \emph{weakly mixing} if every $f \in L^2(\Ebb \mid \Dbb)$
  with $\Exp(f \mid \Dbb) = 0$ is conditionally weakly mixing.
\end{definition}

\begin{remark}[State of the formalization of Lemma~\ref{lem:atomless}]
  The inputs are all in place: the spectral measure exists
  (\texttt{hasSpectralMeasure\_specMeasure}), Wiener's lemma is proved in the form
  \texttt{noAtoms\_of\_tendsto\_cesaro\_zero}, and \eqref{eq:joining} holds for bounded
  functions (\texttt{integral\_joining}).  The Koopman operator is built as a unitary of
  $L^2(X;\CC)$ (\texttt{koopman}), with its powers acting along the orbit
  (\texttt{coeFn\_koopman\_zpow}), and for a bounded real $f$ the correlations
  $\langle f, U^n f\rangle$ are identified with $\int \cip{T^n f}{f} \;dm$
  (\texttt{inner\_toLpC}, \texttt{integral\_eq\_integral\_condInner}).

  \textbf{Part (i) is proved}, as \texttt{noAtoms\_specMeasure\_of\_condWeaklyMixing}.  The step
  from the \texttt{eLpNorm} in the definition of \texttt{CondWeaklyMixing} to the bound
  $\avg{N}\int \abs{\cip{T^n f}{f}} \to 0$ is
  \texttt{tendsto\_cesaro\_integral\_abs\_condInner}, and uses only that $L^1 \leq L^2$ on a
  probability space.  The Fourier coefficients are then dominated by those integrals and by
  $\norm{f}_2 ^2$, so their squares are dominated by a constant times them, and Wiener's lemma
  applies.

  \textbf{Part (ii) is proved} as well, as
  \texttt{noAtoms\_specMeasure\_prod\_of\_condWeaklyMixing}.  The product system
  $(X\times X,\lambda,T\times T)$ is \texttt{prodSystem}, its measure preservation being
  \texttt{joining\_invariant}; the identity \eqref{eq:psi-identity} is
  \texttt{integral\_tensorSq\_comp}, which is \eqref{eq:joining} applied to
  $a = a' = f \cdot (f\circ T^n)$.  The bound $\psi(n) \leq C^2 \phi(n)$ needs
  $\abs{\cip{T^nf}{f}} \leq C ^2$ pointwise, which is \texttt{abs\_condExp\_le}, proved from
  \texttt{condExp\_mono} against the constants $\pm C^2$.
\end{remark}

\begin{definition}[Conditionally compact function; compact extension]
  \label{def:compact-extension}
  \lean{CantorSzemeredi.MPS.CondCompact, CantorSzemeredi.MPS.CompactExtension,
    CantorSzemeredi.MPS.MemModule}
  \leanok
  \uses{def:hilbert-module}
  $f \in L^2(\Ebb \mid \Dbb)$ is \emph{conditionally compact} if for every $\epsilon > 0$ there
  are finitely many $g_1,\dots,g_r \in L^2(\Ebb \mid \Dbb)$ with
  \[
    \sup_{n \in \ZZ} \ \min_{1 \leq q \leq r} \ \bigl\lVert \cnorm{T^n f - g_q} \bigr\rVert_{L^\infty(\Dbb)} < \epsilon .
  \]
  The extension $\Ebb \mid \Dbb$ is \emph{compact} if the conditionally compact functions are
  dense in $L^2(\Ebb)$.
\end{definition}

\begin{remark}
  These definitions are those of Moreira's notes, except that conditional weak mixing is asked
  only for $g = f$ and with ordinary Ces\`aro means.  This is a \emph{weaker} requirement on
  $f$, hence a \emph{weaker} hypothesis on the extension in
  Definition~\ref{def:cwm} than Moreira's --- which is the direction that makes
  Lemma~\ref{lem:dichotomy} apply as stated.  This asymmetry must be reproduced faithfully in
  Lean or the dichotomy will not typecheck against the extension principles.
\end{remark}

\section{The relatively independent self-joining}

Since $X$ is a compact metric space, $m$ disintegrates over $\Dbb$ as $m = \int m_x \, dm(x)$,
so that $\Exp(a \mid \Dbb)(x) = \int a \, dm_x$ for $a \in L^\infty(X)$ and a.e.\ $x$.

\begin{definition}[Relatively independent self-joining]
  \label{def:joining}
  \lean{CantorSzemeredi.MPS.joining, CantorSzemeredi.MPS.condKernel}
  \leanok
  \uses{def:factor}
  $\lambda \coloneqq \int m_x \otimes m_x \, dm(x)$, a probability measure on $X \times X$.
\end{definition}

\begin{lemma}[Properties of $\lambda$]
  \label{lem:joining}
  \lean{CantorSzemeredi.MPS.joining_invariant, CantorSzemeredi.MPS.joining_prod_apply,
    CantorSzemeredi.MPS.condKernel_comp, CantorSzemeredi.MPS.prodMap}
  \uses{def:joining, lem:condexp-commutes}
  For $a, a' \in L^\infty(X)$,
  \begin{equation}
    \label{eq:joining}
    \int_{X \times X} a \otimes a' \, d\lambda = \int_X \Exp(a \mid \Dbb)\, \Exp(a' \mid \Dbb) \, dm ,
  \end{equation}
  and $\lambda$ is invariant under $R \coloneqq T \times T$.
\end{lemma}

\begin{proof}
  \uses{def:joining, lem:condexp-commutes}
  \eqref{eq:joining} is immediate from the disintegration.  For invariance, by
  Lemma~\ref{lem:condexp-commutes} and \eqref{eq:joining},
  $\lambda(R^{-1}(B \times B')) = \int (\Exp(\one_B \mid \Dbb)\Exp(\one_{B'} \mid \Dbb)) \circ T \, dm
   = \lambda(B \times B')$; rectangles generate and are closed under intersection.
\end{proof}

\begin{remark}[State of the formalization]
  \leanok
  The invariance of $\lambda$ under $R = T\times T$ is proved in Lean
  (\texttt{joining\_invariant}), via the rectangle formula
  $\lambda(B\times B') = \int \kappa_x(B)\,\kappa_x(B')\;dm$
  (\texttt{joining\_prod\_apply}), the equivariance $\kappa_x(T^{-1}B) = \kappa_{Tx}(B)$
  (\texttt{condKernel\_comp}, which is Lemma~\ref{lem:condexp-commutes} read through the
  kernel), and the fact that measurable rectangles form a $\pi$-system generating the product
  $\sigma$-algebra.

  The disintegration itself is Mathlib's \texttt{ProbabilityTheory.condExpKernel}, which exists
  for standard Borel spaces --- exactly the standing assumption that $X$ is compact metric with
  its Borel $\sigma$-algebra.

  What is \emph{not} yet formalized is \eqref{eq:joining} for general bounded $a,a'$, as opposed
  to indicators: that needs the Bochner integral against a measure composed with a kernel, which
  Mathlib does not currently provide, and so requires the usual approximation by simple
  functions.  It is used only in Lemma~\ref{lem:atomless}, which is in any case blocked on the
  spectral measure (Definition~\ref{def:spectral-measure}).
\end{remark}

\begin{lemma}[Conditional weak mixing gives atomless spectral measures]
  \label{lem:atomless}
  \lean{CantorSzemeredi.MPS.koopman, CantorSzemeredi.MPS.coeFn_koopman_zpow,
    CantorSzemeredi.MPS.inner_koopman_zpow_eq_integral, CantorSzemeredi.MPS.toLpC,
    CantorSzemeredi.MPS.inner_toLpC, CantorSzemeredi.MPS.integral_eq_integral_condInner,
    CantorSzemeredi.MPS.tendsto_cesaro_integral_abs_condInner,
    CantorSzemeredi.MPS.noAtoms_specMeasure_of_condWeaklyMixing,
    CantorSzemeredi.MPS.prodSystem, CantorSzemeredi.MPS.integral_tensorSq_comp,
    CantorSzemeredi.MPS.noAtoms_specMeasure_prod_of_condWeaklyMixing}
  \leanok
  \uses{def:cwm, lem:joining, lem:wiener, def:spectral-measure}
  Let $f \in L^\infty(\Ebb)$ be conditionally weakly mixing and set $G \coloneqq f \otimes \bar f$.
  Then
  \begin{enumerate}
    \item[(i)] $\sigma_f$, for the Koopman operator of $T$ on $L^2(X)$, has no atoms;
    \item[(ii)] $\sigma_G$, for the Koopman operator of $R$ on $L^2(\lambda)$, has no atoms, and
      \begin{equation}
        \label{eq:psi-identity}
        \langle R^n G, G \rangle_{L^2(\lambda)} = \bigl\lVert \cip{T^n f}{f} \bigr\rVert_{L^2(\Dbb)}^2 ,
        \qquad n \in \ZZ .
      \end{equation}
  \end{enumerate}
\end{lemma}

\begin{proof}
  \uses{def:cwm, lem:joining, lem:wiener}
  \eqref{eq:psi-identity} is \eqref{eq:joining} with $a = T^n f \cdot \bar f$ and $a' = \bar a$.
  Normalize $\norm{f}_\infty \leq 1$ and put $a_n \coloneqq \lVert \cip{T^n f}{f} \rVert_{L^1(\Dbb)}$.
  Since the $\abs{\cip{T^n f}{f}}$ are non-negative and bounded by $1$,
  $\avg{1 \leq n \leq N} a_n \leq \lVert \avg{1 \leq n \leq N} \abs{\cip{T^n f}{f}} \rVert_{L^2(\Dbb)} \to 0$.
  Now $\abs{\widehat{\sigma_f}(n)} \leq a_n$ and $0 \leq \widehat{\sigma_G}(n) \leq a_n$ by
  \eqref{eq:psi-identity}.  As $a_n \leq 1$,
  $\avg{1 \leq n \leq N} \abs{\widehat\sigma(n)}^2 \leq \avg{1 \leq n \leq N} a_n \to 0$ for
  $\sigma \in \{\sigma_f, \sigma_G\}$, and Wiener's lemma (Lemma~\ref{lem:wiener}) applies.
\end{proof}

\section{Passing to powers}

\begin{lemma}[Powers]
  \label{lem:powers}
  \lean{CantorSzemeredi.MPS.CompactExtension.pow, CantorSzemeredi.MPS.CondCompact.pow,
    CantorSzemeredi.MPS.pow, CantorSzemeredi.MPS.pow_iter_eq}
  \uses{def:compact-extension, def:cwm}
  Let $k \geq 1$.
  \begin{enumerate}
    \item If $\Ebb$ is a compact extension of $\Dbb$ for $T$, then it is one for $T^k$.
    \item If $f, g \in L^2(\Ebb \mid \Dbb)$ satisfy
      $\lim_N \lVert \avg{N} \abs{\cip{T^n f}{g}} \rVert_{L^2(\Dbb)} = 0$, then the same holds
      along $k\NN$.  In particular, if hypothesis~\eqref{eq:WM} holds for $T$ it holds for $T^k$.
  \end{enumerate}
\end{lemma}

\begin{proof}
  \uses{def:compact-extension, def:cwm}
  (1) $\{T^{kn} f\} \subset \{T^n f\}$, so conditional compactness is inherited and the
  conditionally compact functions remain dense.

  (2) Set $\phi_n \coloneqq \abs{\cip{T^n f}{g}} \geq 0$.  Since
  $\{kn \colon n<N\} \subset \{m \colon m<kN\}$ and the $\phi_m$ are non-negative, the
  subsampled average is dominated pointwise by $k$ times the full average,
  \begin{equation} \label{eq:subsample}
    \avg{n<N} \phi _{kn} \leq k \cdot \avg{m<kN} \phi _m .
  \end{equation}
  Taking $L^2(\Dbb)$ norms and letting $N \to \infty$ gives the claim.
\end{proof}

\begin{remark}[A simplification found while formalizing]
  \leanok
  The proof of (2) just given is shorter than the one in the first draft of the paper, which
  passed through the characterization of $\avg{N}a_n \to 0$ by convergence of $a_n$ along a set
  of density one, and then argued that a $k$-subsampling of a density one set again has density
  one.  That detour is unnecessary: \eqref{eq:subsample} is a pointwise inequality between
  non-negative quantities, and costs only the harmless constant $k$.  No density argument, no
  Cauchy--Schwarz, and no $L^1$/$L^2$ comparison is needed.  In Lean the inequality
  \eqref{eq:subsample} is \texttt{CantorSzemeredi.MPS.avg\_subsample\_le}.
\end{remark}

\section{The van der Corput inequality}

\begin{lemma}[van der Corput]
  \label{lem:vdc}
  \lean{CantorSzemeredi.VdCData, CantorSzemeredi.corr,
    CantorSzemeredi.VdCData.van_der_corput, CantorSzemeredi.VdCData.van_der_corput_avg}
  \leanok
  Let $u_M,\dots,u_{N-1}$ be vectors of norm at most $1$ in a Hilbert space and $1 \leq H < N-M$.
  Then
  \begin{equation}
    \label{eq:vdc}
    \Bigl\lVert \unifavg u_n \Bigr\rVert^2
      \leq 4 \avg{1 \leq h \leq H} \Bigl\lvert \avg{M \leq n < N-h} \langle u_{n+h}, u_n \rangle \Bigr\rvert
        + \frac{2}{H} .
  \end{equation}
\end{lemma}

\begin{proof}
  \leanok
  Write $L = N-M$ and set $u_n = 0$ off $[M,N)$.  Each $u_n$, $n \in [M,N)$, occurs exactly $H$
  times among the $u_{m+h}$ with $1 \leq h \leq H$ and $m \in I \coloneqq [M-H, N-1)$, an
  interval of fewer than $L+H$ integers.  Cauchy--Schwarz in $m$ gives
  $\lVert \sum_n u_n \rVert^2 \leq \frac{L+H}{H^2}\sum_{h,h'=1}^{H}\sum_{m \in I}\langle u_{m+h}, u_{m+h'}\rangle$.
  The $H$ diagonal terms contribute at most $L$ each.  For $r = h-h' \neq 0$ the inner sum is
  $S(r) = \sum_n \langle u_{n+r}, u_n\rangle$ with $S(-r) = \overline{S(r)}$, and exactly
  $H - \abs{r}$ pairs give $h - h' = r$; moreover $\abs{S(r)} \leq L\abs{c_r}$ with
  $c_r = \avg{M \leq n < N-r}\langle u_{n+r}, u_n\rangle$.  Combining,
  $\lVert \unifavg u_n \rVert^2 \leq \frac1H + \frac1L + \frac{2(L+H)}{L}\avg{1 \leq r \leq H}\abs{c_r}$,
  and $H < L$ gives \eqref{eq:vdc}.
\end{proof}

\begin{remark}[Lean form]
  \leanok
  The Lean statement \texttt{van\_der\_corput} is slightly stronger: its right hand side is
  $4\, \avg{1 \leq h \leq H} \abs{S(h)}/L + 2/H$, with $S(h) = \sum_{M \leq n<N} \langle u_{n+h}, u_n\rangle$
  divided by the full length $L$, rather than $\abs{c_h} = \abs{S(h)}/(L-h)$.  Since $L-h \leq L$ this
  implies \eqref{eq:vdc}, which is \texttt{van\_der\_corput\_avg}.

  Two points of the formalization are worth recording.  Indices run over $\ZZ$ rather than
  $\NN$: the proof sums over $m \in [M-H,N)$, and with $\NN$ subtraction $M-H$ truncates at $0$
  when $M<H$, which breaks the covering count.  And the vectors are taken in a \emph{real}
  inner product space, which makes $S(-r)=S(r)$ the symmetry of the real inner product rather
  than a statement about conjugates; everything this lemma is applied to is real.
\end{remark}
